\section{Paganism 2.0}

\begin{quotex}
I see and approve what is better, but follow what is worse.
\flright{\textsc{Ovid}, \emph{Metamorphoses}}

I do not do the good I want, but the evil I do not want is what I do.
\flright{\textsc{Paul}, \emph{Romans 7}}
\end{quotex}


For some reason, there are some who think it is a simple matter to create a religion \emph{de novo}, often with the intent to trick a population into pursuing some desired end. Historically, however, starting with an existing, respected religion and adding onto it seems to have a better chance of success. Recent examples include Joseph Smith with Mormonism and Mary Baker Eddy with Christian Science.

For a variety of reasons which are not important for our purpose, there are occasional calls for a return to a pre-Christian paganism. This is often tied in with the belief in a more ``natural'', almost Edenic, lifestyle, unpolluted with the Christian consciousness of the loss of original innocence. Unfortunately, that state never existed, as our epigraph shows. Rather, the pagan was also in torment, feeling himself tied to a fate beyond his control, or at the mercy of boons and banes coming from the gods.

This is not to deny that a noble attitude could not arise, far from it. But the modern neo-pagan does not believe in that in the same way. Moreover, the ``mule'' of neo-paganism is burdened with things few want to return to: raids and conquests for personal gain, human sacrifice, slavery, a real metaphysical pluralism, and so on. A new paganism, therefore, should build on what currently exists and try to move beyond it.

We can point to two attempts in the 20\textsuperscript{th} century at a new pagan philosophy: \textbf{Julius Evola} and \textbf{Charles Maurras}. Although much of what they believed is anathema to the modern world, their systems are worth investigating, if only for historical interest. At a minimum, they set standards for intellectual rigour.

\paragraph{Magical Idealism}


Evola rejected the neopagan project. He even asserted that the neopagans would be better off joining the Church than pursuing an impossible reconstruction. Rather than back-tracking, he proposed, it seems to me, a post-Christian paganism that wants to recapture a certain nobility of spirit without resuscitating the dead mule of paganism.

For example, Paganism 2.0, in his view, would have to incorporate the spirituality of the Middle Ages by separating its pagan elements from the strictly Christian elements. This brings to mind Psyche's task of separating the various grains from the pile. In other words, it requires a super-human effort.

Now Evola took the concept of ``person'' as central, even writing more than once of the ``dignity of the person''. The Person, moreover, is a Christian concept, not part of the pagan mindset. Evola also rejected polytheism, since a philosophical pluralism is non-Traditional. Nevertheless, his idea of God seems to be a ``god in evolution''. That is, God exists to the extent that the Self achieves the Absolute Self. So rather than pluralism, Evola actually defends solipsism. Guenon regards a belief in an evolving god as ``Satanic''.

So what are his raw materials? First of all, there is his philosophical system of magical idealism as expounded in three books (\textit{Phenomenology of Magical Idealism}, \textit{Theory of Magical Idealism}, \textit{Essays in Magical Idealism}) plus the \textit{Yoga of Power}, that is, Tantra Yoga as mediated by John Woodroffe. Magical idealism draws on German and Italian idealism, with insights from Otto Weininger, Max Stirner, Carlo Michelstaedter, Friedrich Nietzsche, Rene Guenon, of course, and various Eastern and Western occult, esoteric, and spiritual traditions. It certainly is an interesting melange.

Beyond that, Evola wrote books on Hermetism and Buddhism. Although the history of Hermetism is entwined with Christianity, as a post-Christian pagan, Evola attempted to extract its aspects that suited him. He also proposed a form of Buddhism, which I don't believe actually exists anywhere.

This seems to have been a purely intellectual exercise, since, apart from some experiments in the Ur/Krur group, there was no actual esoteric path for the actualization of magical idealism. Since Evola was not fond of priests, there was no basis for an esoteric chain. The post-Christian pagan society would be ruled by a class of nobles and warriors, with the priestly role relegated to secondary importance. Perhaps they would be suitable for the merchants and serfs, but the nobles and warriors would follow the ethos of magical idealism, even if they couldn't understand it.

\paragraph{Positivism}


The French Revolution shook up Europe at its foundations. The old regime, based on religion and aristocracy, holdovers from the Middle Ages, was suddenly replaced by a revolution from below. That revolution overturned the established order of things. In its aftermath, the distinction between the Left and the Right arose: the Left was in a state of permanent revolution, while the right was nostalgic for throne and altar. The best representatives of that tendency are \textbf{Edmund Burke} and \textbf{Joseph de Maistre}. Anyone interested in the actual ``\textbf{Old Right}'' needs to start with them.

From that experience, there arose the mathematician and philosopher, \textbf{Auguste Comte}. If the old regime could no longer support the established order of things based on religious ideals, Comte looked for that support instead in science. His philosophy of \textbf{Positivism} regarded as real only that which could be experienced. Hence, he created a hierarchy of the sciences, based on the object of experience. He regarded the soul faculties of thinking, feeling, and willing as experiences, thereby avoiding the scientific reductionism of the positivists who came after him.

Because his system was conservative, and had space in it for the ``higher'' impulses in man that religion had previously claimed for itself, his ideas gained some traction. In particular, he had influence in South America because his social policy was compatible with Catholic thought. Brazil's flag includes one of his mottoes (\emph{Order and Progress}), although I haven't met a Brazilian who knows where it came from.

Charles Maurras was one such thinker who considered Comte as his mentor. He also looked to Catholic social theory, as it was then being formulated by the popes starting with Pope Leo XIII ... that was the original ``third way'' between socialism and capitalism, although neglected today. His other influence was \textbf{Ernest Renan}, the historian who gave up the priesthood and wrote the famous history of Jesus instead. Although he wrote it as a secular history, he wasn't destructive, since he still recognized the ideals of religion while denying its supernatural content. Maurras also took much more from Renan, such as his ideas on nationalism. Nationalists today should follow that cue, rather than the enraged irrationalities coming from neopagans today. Nevertheless, such ideas will be pretty indigestible in Western nations today.

In France, Maurras brought \textbf{Action Française} to its peak of influence and popularity. It brought together secularists and Catholics, since they could agree on social issues. Catholic teaching is that such doctrines are knowable by natural reason, so there is a legitimate reason for cooperation. Even thinkers like \textbf{Jacques Maritain} and \textbf{Rene Guenon} were supporters of Action Française, until it was condemned by Pope Pius XI. When Pius XII lifted the condemnation, the damage had been done and it never recovered its former prestige. Atheistic secularism was given free rein in France.

Neopagan, nationalist rightists would do well to look to these figures for their intellectuality and emotional detachment, since they actually had a following. There is a reason that neopagans are not taken seriously today.


\flrightit{Posted on 2015-05-03 by Cologero}

\begin{center}* * *\end{center}

\begin{footnotesize}
\begin{sffamily}

\texttt{Max on 2015-05-04 at 10:30 said:}

One of the problems with Evola is that despite all talk of asserting masculine values, he sometimes approaches a kind of feminine solipsism. It seems that a true masculinity would have to acknowledge the feminine as well, internally and externally, for creation is not complete without it. Grounding ourselves in positivism, we have to. Otherwise we commit a similar fault as regarding the world as evil, but on a microcosmic order. Drawing from previous posts, it seems the way to overcome a ``privation'' is to love it.

\hfill

\texttt{obcure on 2015-05-04 at 19:29 said:}

Max,

To overcome privations one must simply turn to the Good. Wherever there is any degree of reality there is the Good. Where there is not the Good there is not just privation but negation and thus nothing.

\hfill

\end{sffamily}
\end{footnotesize}

